% 出席確認クイズ 第03回　生成AIのしくみと限界
%   attendance/attendance03.html から生成（解説は本ガイド用に加筆）。

\quizq{1}{大規模言語モデル(LLM)が文章を生成する基本的なしくみはどれでしょうか?}%
  {辞書の意味を順番に並べる}%
  {人手で記述した文法規則をすべて組み合わせる}%
  {人が書いた文章をデータベースからそのまま検索して表示する}%
  {\correct{大量のテキストから学習し、次に来る語(トークン)の確率を予測して生成する}}%
  {正解は (d)。大規模言語モデル（LLM）は膨大なテキストから学習したパターンをもとに、次に来るトークンの確率を予測しながら文章を1語ずつ生成します。(c) のようにデータベースを検索して表示しているわけではないため、出典が存在しない内容でも流暢に生成してしまいます。これがハルシネーションの根本的な原因です。(a)(b) は古典的な手法の説明です。}

\quizq{2}{「ハルシネーション」の説明として適切なものはどれでしょうか?}%
  {AIの動作が遅くなること}%
  {\correct{AIが事実に基づかない内容をもっともらしく生成すること}}%
  {AIが同じ回答を繰り返すこと}%
  {AIが画像を認識できないこと}%
  {正解は (b)。ハルシネーションは、事実に基づかない内容をもっともらしく、自信のある口調で生成してしまう現象です。(a)(c)(d) は別の問題です。実在する資料名に存在しない数値を添えるなど「半分正しい」形で現れることが多く、固有名詞と数値は必ず一次資料で確かめる必要があります。}

\quizq{3}{大規模言語モデルが一般に苦手なこととして適切なものはどれでしょうか?}%
  {言語の翻訳}%
  {文章の要約}%
  {文体の変換}%
  {\correct{学習データ以降の最新情報を正確に答えること}}%
  {正解は (d)。LLM の知識には学習データの締め切り（知識カットオフ）があり、それ以降の出来事は知りません。Web検索や社内文書の検索と組み合わせる RAG（検索拡張生成）などで補います。(a)(b)(c) の翻訳・要約・文体変換は、LLM の得意分野です。}

\quizq{4}{「プロンプト」の説明として適切なものはどれでしょうか?}%
  {\correct{AIに与える指示や入力文}}%
  {AIの学習に使うデータセット}%
  {AIの処理速度の単位}%
  {AIが出力する文章}%
  {正解は (a)。プロンプトは AI に与える指示や質問の文です。(b) は学習データ、(d) は出力（応答）にあたります。本講義では、役割・目的と背景・条件・出力の形式という4要素を意識してプロンプトを書きます。}

\quizq{5}{「トークン」の説明として適切なものはどれでしょうか?}%
  {セキュリティのための合言葉}%
  {画像の解像度の単位}%
  {\correct{モデルが扱う文章の最小単位(語や部分語)}}%
  {AIの利用料金のこと}%
  {正解は (c)。トークンは LLM が文章を処理する最小単位で、単語や単語の一部に相当します。日本語では1文字が1〜2トークンになることが多いです。文脈窓（一度に扱える長さ）や API の料金はトークン数で決まりますが、(d) のようにトークンそのものが料金なのではありません。}
